\documentclass{abstract}

% Bibliography
\usepackage[
    backend=biber,
    style=authoryear,
    doi=true,
    sorting=none,
    sortcites=true,
    maxbibnames=2,
    minbibnames=1,
    maxcitenames=2,
    mincitenames=1,
    citestyle=numeric-comp,
    giveninits=true,
    isbn=false,
    date=year
]{biblatex}
\addbibresource{abstract.bib}

% For email linking
\usepackage{hyperref}

\title{TITLE HERE}

\author{Jessica Zhang}

\keywords{cone calorimeter, database, flammability, SmURF}
\begin{document}
\maketitle
\section{Abstract}
The cone calorimeter is a tool that used to evaluate different aspects of flammable materials, including heat release rate. However, there has been no centralized database to store and organize decades of cone calorimeter data in an accessible format. This is essential for fire safety research, as understanding the flammability of materials can provide insight into how fire disasters occur and help improve fire protection strategies and standards. To address this need, a methodology called SmURF (Systematic User Review of Files) was developed to systematically review published cone calorimeter reports, correct PDF parsing errors, assign standardized material IDs, and prepare data to be incorporated into the larger Material Flammability database. This project primarily involves SmURFing multiple reports and using an AI tool to reparse PDFs to improve the accuracy of previously unsuccessful parsing attempts.

Example citation \cite{babrauskas1982development, brown1988cone,emil1989flammability}


\section{References}
\printbibliography[heading=none]

\end{document}